% GENERATED FILE. Edit canonical lessons, then run npm run generate:books.
% canonical-source-hash: fnv1a64:9943febc27906111
% canonical-lessons: TE-C115-topi, TE-C115-ungaram, TE-C115-gajulu, TE-C115-kura, TE-C115-pappu

\chapter{Cap, Ring, Bangles, Curry, Dal}
\label{ch:te-cap-ring-bangles-curry-dal}
\hlchaptermodality{115}

\begin{chapteropening}
\textbf{By the end of this chapter:} \emph{I can say a cap, a ring, bangles, a curry and dal.}

Complete the last lesson of chapter 115: i can say a cap, a ring, bangles, a curry and dal.
\end{chapteropening}

\section[ṭōpī]{\te{టోపీ} (ṭōpī) — a cap, a hat}

\label{lesson:TE-C115-topi}



\begin{quote}
Before the new one: say the Telugu for a dhoti, then the Telugu for sandals.
\end{quote}

\subsection*{You'll want to know: \te{టోపీ}}
\textbf{\te{టోపీ}} — \emph{ṭōpī} — \textquotedblleft{}a cap, a hat\textquotedblright{}.

\begin{grammarlens}[title={What you've built}]
One more word for this chapter.
\end{grammarlens}

\subsection*{Your turn}
\begin{itemize}
\raggedright
  \item \emph{Say it:} \emph{ṭōpī}
  \item \emph{Say it:} \emph{ṭōpī}, once more
  \item \emph{Say it:} say \emph{ṭōpī}
  \item \emph{Recall:} say the Telugu for a dhoti, then the Telugu for sandals, then say \emph{ṭōpī} again
\end{itemize}

\begin{tcolorbox}[breakable,colback=teal!4,colframe=teal!35!black,title={Before you move on}]
What does \te{టోపీ} mean? (\textquotedblleft{}A cap, a hat\textquotedblright{}.) Say it once more.
\end{tcolorbox}
\section[uṅgaraṁ]{\te{ఉంగరం} (uṅgaraṁ) — a ring}

\label{lesson:TE-C115-ungaram}



\begin{quote}
Before the new one: say the Telugu for sandals, then the Telugu for a cap.
\end{quote}

\subsection*{You'll want to know: \te{ఉంగరం}}
\textbf{\te{ఉంగరం}} — \emph{uṅgaraṁ} — \textquotedblleft{}a ring\textquotedblright{}.

\begin{grammarlens}[title={What you've built}]
One more word for this chapter.
\end{grammarlens}

\subsection*{Your turn}
\begin{itemize}
\raggedright
  \item \emph{Say it:} \emph{uṅgaraṁ}
  \item \emph{Say it:} \emph{uṅgaraṁ}, once more
  \item \emph{Say it:} say \emph{uṅgaraṁ}
  \item \emph{Recall:} say the Telugu for sandals, then the Telugu for a cap, then say \emph{uṅgaraṁ} again
\end{itemize}

\begin{tcolorbox}[breakable,colback=teal!4,colframe=teal!35!black,title={Before you move on}]
What does \te{ఉంగరం} mean? (\textquotedblleft{}A ring\textquotedblright{}.) Say it once more.
\end{tcolorbox}
\section[gājulu]{\te{గాజులు} (gājulu) — bangles}

\label{lesson:TE-C115-gajulu}



\begin{quote}
Before the new one: say the Telugu for a cap, then the Telugu for a ring.
\end{quote}

\subsection*{You'll want to know: \te{గాజులు}}
\textbf{\te{గాజులు}} — \emph{gājulu} — \textquotedblleft{}bangles\textquotedblright{}.

\begin{grammarlens}[title={What you've built}]
One more word for this chapter.
\end{grammarlens}

\subsection*{Your turn}
\begin{itemize}
\raggedright
  \item \emph{Say it:} \emph{gājulu}
  \item \emph{Say it:} \emph{gājulu}, once more
  \item \emph{Say it:} say \emph{gājulu}
  \item \emph{Recall:} say the Telugu for a cap, then the Telugu for a ring, then say \emph{gājulu} again
\end{itemize}

\begin{tcolorbox}[breakable,colback=teal!4,colframe=teal!35!black,title={Before you move on}]
What does \te{గాజులు} mean? (\textquotedblleft{}Bangles\textquotedblright{}.) Say it once more.
\end{tcolorbox}
\section[kūra]{\te{కూర} (kūra) — a curry}

\label{lesson:TE-C115-kura}



\begin{quote}
Before the new one: say the Telugu for a ring, then the Telugu for bangles.
\end{quote}

\subsection*{You'll want to know: \te{కూర}}
\textbf{\te{కూర}} — \emph{kūra} — \textquotedblleft{}a curry\textquotedblright{}.

\begin{grammarlens}[title={What you've built}]
One more word for this chapter.
\end{grammarlens}

\subsection*{Your turn}
\begin{itemize}
\raggedright
  \item \emph{Say it:} \emph{kūra}
  \item \emph{Say it:} \emph{kūra}, once more
  \item \emph{Say it:} say \emph{kūra}
  \item \emph{Recall:} say the Telugu for a ring, then the Telugu for bangles, then say \emph{kūra} again
\end{itemize}

\begin{tcolorbox}[breakable,colback=teal!4,colframe=teal!35!black,title={Before you move on}]
What does \te{కూర} mean? (\textquotedblleft{}A curry\textquotedblright{}.) Say it once more.
\end{tcolorbox}
\section[pappu]{\te{పప్పు} (pappu) — dal, lentils}

\label{lesson:TE-C115-pappu}



\begin{quote}
Before the new one: say the Telugu for bangles, then the Telugu for a curry.
\end{quote}

\subsection*{You'll want to know: \te{పప్పు}}
\textbf{\te{పప్పు}} — \emph{pappu} — \textquotedblleft{}dal, lentils\textquotedblright{}.

\begin{grammarlens}[title={What you've built}]
One more word for this chapter.
\end{grammarlens}

\subsection*{Your turn}
\begin{itemize}
\raggedright
  \item \emph{Say it:} \emph{pappu}
  \item \emph{Say it:} \emph{pappu}, once more
  \item \emph{Say it:} say \emph{pappu}
  \item \emph{Recall:} say the Telugu for bangles, then the Telugu for a curry, then say \emph{pappu} again
\end{itemize}

\begin{tcolorbox}[breakable,colback=teal!4,colframe=teal!35!black,title={Before you move on}]
What does \te{పప్పు} mean? (\textquotedblleft{}Dal, lentils\textquotedblright{}.) Say it once more.
\end{tcolorbox}
